\section{Évaluation des systèmes RAG}

\subsection{Séparer les composants évalués}

L'évaluation d'un assistant RAG ne peut pas être réduite à une seule mesure. Une réponse insatisfaisante peut résulter d'un document absent, d'un mauvais découpage, d'une récupération incorrecte, d'une mauvaise interprétation ou d'un problème d'interface. Les travaux RAGAS proposent notamment de distinguer la pertinence du contexte, la fidélité au contexte et la pertinence de la réponse \citep{Es2024RAGAS}. ARES propose également une évaluation automatisée de la pertinence du contexte, de la fidélité et de la pertinence de la réponse, accompagnée d'un nombre limité d'annotations humaines \citep{SaadFalcon2024ARES}.

Ces cadres sont utiles pour structurer un protocole, mais les juges automatiques peuvent eux-mêmes se tromper. Une évaluation crédible combine des métriques automatiques, des vérifications humaines et des tests techniques reproductibles.

\begin{table}[H]
\centering
\caption{Dimensions d'évaluation applicables à l'Assistant RAG DNSI}
\label{tab:evaluation_dimensions}
\begin{tabular}{p{3.2cm}p{6cm}p{4.5cm}}
\toprule
\textbf{Dimension} & \textbf{Question évaluée} & \textbf{Exemples de preuves} \\
\midrule
Ingestion & Les documents attendus sont-ils extraits, découpés, indexés et mis à jour ? & États de synchronisation, comptages, tests de suppression et de mise à jour. \\
Retrieval & Les fragments pertinents apparaissent-ils parmi les premiers résultats ? & Recall@$k$, précision@$k$, jugement de pertinence sur un jeu de questions. \\
Génération & La réponse est-elle fidèle aux passages et répond-elle à la question ? & Exactitude factuelle, fidélité, citations, analyse humaine. \\
Sécurité & Un utilisateur ne récupère-t-il que les documents autorisés ? & Tests positifs et négatifs par profils, cas d'ACL manquantes et héritées. \\
Système & Le service fonctionne-t-il de manière stable sous la charge attendue ? & Latence par composant, erreurs, saturation GPU, disponibilité observée. \\
Usage & La réponse et les sources sont-elles compréhensibles et utiles ? & Scénarios utilisateurs, feedbacks, entretiens ou questionnaire si réellement réalisés. \\
\bottomrule
\end{tabular}
\end{table}

\subsection{Traçabilité et reproductibilité des résultats}

Une mesure n'est exploitable dans un mémoire que si son origine, son protocole de collecte et ses conditions d'obtention peuvent être établis. Les résultats retenus dans le chapitre d'évaluation devront donc pouvoir être reliés à un script, un journal d'exécution, un jeu de données, un export ou un protocole reproductible.

Lorsque certaines dimensions n'ont pas fait l'objet d'une mesure quantitative complète, l'analyse s'appuiera sur les tests fonctionnels disponibles, les observations vérifiables et la description d'un protocole d'évaluation complémentaire. Cette distinction entre résultats observés et évaluations à conduire permet de conserver une interprétation rigoureuse des performances de la solution.

\section{Positionnement de l'Assistant RAG DNSI}

\begin{table}[H]
\centering
\caption{Correspondance entre l'état de l'art et la stack réellement observée}
\label{tab:positionnement_stack}
\begin{tabular}{p{3.5cm}p{5.2cm}p{5.2cm}}
\toprule
\textbf{Concept} & \textbf{Choix du projet} & \textbf{Précaution d'interprétation} \\
\midrule
Embeddings denses & \texttt{multilingual-e5-large}, préfixes query/passage, normalisation des vecteurs. & Les performances sur le corpus DNSI doivent être établies par une évaluation locale. \\
Base vectorielle & Qdrant avec payloads, collections, URL et ACL. & Les composants FAISS conservés dans le dépôt correspondent à des chemins historiques ou de compatibilité. \\
Retrieval & Recherche dense puis reranking lexical et heuristique. & Les règles métier doivent être testées et maintenues ; ce mécanisme ne constitue pas un reranker neuronal appris. \\
Génération & Mistral Small 3.1 24B servi localement par vLLM. & Les performances publiques du modèle ne constituent pas des résultats propres au projet. \\
Sources & Connecteurs SharePoint et Freshservice. & Le niveau de couverture doit être apprécié à partir des flux effectivement configurés et testés. \\
Sécurité & Principaux normalisés et filtre ACL Qdrant avant génération. & Le comportement en présence de métadonnées incomplètes doit faire l'objet de tests négatifs. \\
Déploiement & Docker Compose, PostgreSQL, Nginx, healthchecks et VM GPU. & Cette architecture décrit le déploiement observé, sans préjuger d'une haute disponibilité multi-nœuds. \\
\bottomrule
\end{tabular}
\end{table}

Le système est ainsi positionné comme une application RAG d'entreprise modulaire, déployée sur une infrastructure interne et connectée à plusieurs sources documentaires. Sa contribution principale réside dans l'intégration cohérente des composants, dans l'adaptation du pipeline aux contenus de la DNSI et dans la prise en compte des contraintes de traçabilité, d'habilitation et d'exploitation.

\section{Conclusion}

L'état de l'art met en évidence la complémentarité des composants mobilisés dans l'Assistant RAG DNSI. Les Transformers fournissent les modèles d'encodage et de génération. Les bi-encodeurs permettent de représenter les questions et les passages dans un espace sémantique. Le RAG associe cette recherche à un modèle génératif, tandis que Qdrant apporte la persistance, les payloads et le filtrage nécessaire aux collections documentaires. vLLM assure le service local du modèle Mistral, et les connecteurs adaptent le pipeline aux contraintes de SharePoint et Freshservice.

La littérature montre également qu'aucun de ces composants ne garantit seul la fiabilité. La qualité dépend de la chaîne complète : collecte, découpage, embeddings, recherche, autorisations, contexte, génération, affichage des sources et évaluation. Le chapitre suivant présente la démarche méthodologique employée pour concevoir et valider cet artefact en distinguant les objectifs, les observations et les résultats effectivement établis.
